\documentclass[11pt]{article}

\usepackage[margin=1in]{geometry}
\usepackage{amsmath, amssymb}
\usepackage{booktabs}
\usepackage{hyperref}
\usepackage{siunitx}

\title{Regime-Switching Market Making: \\ Filter and Optimal Control Corrections}
\author{Session notes}
\date{\today}

\newcommand{\E}{\mathbb{E}}
\newcommand{\Var}{\mathrm{Var}}
\newcommand{\N}{\mathcal{N}}

\begin{document}
\maketitle

\begin{abstract}
This note documents, from a theoretical standpoint, a set of implementation errors identified and corrected in the Hamilton-filter regime-detection module and the associated optimal-control comparison (oracle / belief-weighted / randomised / naive agents) built on top of \texttt{mbt\_gym}. Each error is presented as a mismatch between an assumed model and the model actually implemented in code, together with the correction and, where useful, an empirical measurement of its impact. The note closes with a discussion of why perfect regime information does not, by itself, guarantee the oracle agent dominates on raw mean PnL.
\end{abstract}

\tableofcontents

\section{Setup and notation}

The environment simulates a single-asset limit order book market maker operating under a hidden two-state regime $Z_t \in \{0,1\}$, evolving as a discrete-time Markov chain with transition matrix $P$. Regime $0$ is calm (pure Brownian diffusion); regime $1$ is adverse selection (diffusion plus a permanent price impact on every market order arrival).

The midprice evolves, in \emph{absolute price units}, as
\begin{equation}
S_{t+\Delta t} = S_t + \sigma_{Z_t} \sqrt{\Delta t}\, \xi_t + \epsilon\big(\mathbb{1}[\text{buy MO arrival}] - \mathbb{1}[\text{sell MO arrival}]\big), \qquad \xi_t \sim \N(0,1),
\label{eq:midprice}
\end{equation}
where the jump term is present only in regime $1$. Market order arrivals on each side are Poisson with intensity $\lambda$, so the probability of an arrival in a step of size $\Delta t$ is
\begin{equation}
p = \lambda \Delta t.
\label{eq:jumpprob}
\end{equation}

An agent quotes bid/ask depths $\delta^{b}, \delta^{a}$ relative to $S_t$; fills occur with probability $\propto e^{-\kappa \delta}$ (exponential fill intensity, decay parameter $\kappa$). The reward accrued is the \emph{running inventory penalty} criterion
\begin{equation}
J = \E\left[\int_0^T dC_t + q_t\, dS_t - \phi\, q_t^2\, dt \right] - \alpha\, q_T^2,
\label{eq:objective}
\end{equation}
where $C_t$ is cash, $q_t$ is inventory, $\phi$ is the running inventory-aversion coefficient and $\alpha$ is the terminal inventory-aversion coefficient.

A \texttt{HamiltonFilter} object tracks the posterior belief $\pi_t = P(Z_t = 1 \mid \mathcal{F}_t)$ from the observed midprice series, using an emission model over the \emph{percentage return}
\begin{equation}
r_t = \frac{S_t - S_{t-\Delta t}}{S_{t-\Delta t}}.
\label{eq:pctreturn}
\end{equation}
A separate closed-form (semi-analytic) optimal control, described in Section~\ref{sec:control}, is used to build an \emph{oracle} agent (knows $Z_t$ exactly) and to interpolate/sample controls for belief-driven agents.

Seven distinct implementation errors were found across this pipeline. Each broke a different theoretical identity between the assumed and implemented model; none were visible from a single stack trace, and each was confirmed empirically before being corrected.

\section{Error 1: Observation normalisation vs.\ raw state}
\label{sec:norm}

\texttt{mbt\_gym}'s \texttt{TradingEnvironment} applies, by default, an affine normalisation to every returned observation so that it lies in $[-1,1]$:
\begin{equation}
\tilde{x} = \frac{x - \ell}{(u - \ell)/2} - 1,
\label{eq:obsnorm}
\end{equation}
where $[\ell, u]$ are the (regime-specific) bounds of the raw observation space for that component. In particular, the cash bound is
\begin{equation}
u_{\text{cash}} = n_{\text{steps}} \cdot \bar{S},
\qquad \bar S = S_0 + \max\text{-price-move for that regime},
\end{equation}
which differs between the calm and adverse-selection sub-environments because their midprice models have different price ranges (the adverse-selection model's range is dominated by the jump term in \eqref{eq:midprice}). Concretely, in this configuration $u_{\text{cash}}^{(0)} \approx 20{,}008$ and $u_{\text{cash}}^{(1)} \approx 100{,}080$ — a factor of $\sim 5$ apart.

\paragraph{Consequence.} Two implementation patterns both treated $\tilde x$ as if it were $x$:
\begin{enumerate}
\item The filter's midprice input used the normalised observation, so the ``percentage return'' \eqref{eq:pctreturn} computed on $\tilde S$ instead of $S$ produced values of order $10^{-1}$ per step instead of the true $\sim 10^{-5}$--$10^{-2}$, i.e.\ the input series bore no resemblance to a real return series.
\item PnL bookkeeping computed $\tilde C_T - \tilde C_0$. Since the normalisation gradient $(u-\ell)/2$ is regime-dependent, this quantity is not even a consistently-scaled proxy for $C_T - C_0$ whenever a regime switch occurs mid-episode (which is most episodes, given the transition dynamics) — it silently mixes two different linear scalings.
\end{enumerate}

\paragraph{Correction.} A \texttt{raw\_state} accessor was added to the environment wrapper, exposing $(C_t, q_t, t, S_t)$ read directly from the underlying \texttt{model\_dynamics.state} (pre-normalisation), captured at the moment each regime is active (before any regime transition is applied for the next step). All downstream cash/inventory/price bookkeeping and filter input now use this raw state.

\section{Error 2: Volatility units (absolute price vs.\ percentage)}
\label{sec:pctunits}

The midprice model \eqref{eq:midprice} uses $\sigma_{Z_t}$ as an \emph{absolute price} diffusion coefficient. The filter, however, is defined over the percentage return \eqref{eq:pctreturn}. For small $\Delta S$,
\begin{equation}
r_t \approx \frac{\Delta S_t}{S_0}, \qquad \Rightarrow \qquad \Var(r_t) \approx \frac{\sigma_{Z_t}^2 \Delta t}{S_0^2}.
\end{equation}
The filter's constructor was passed $\sigma_{Z_t}$ directly (in price units, $\mathcal O(10^{-2}\text{--}10^{-1})$), rather than the percentage-scaled version. Since the true per-step percentage volatility here is $\mathcal O(10^{-6}\text{--}10^{-5})$, the assumed $\sigma$ was of the order $10^3$--$10^4$ times too large.

\paragraph{Correction.}
\begin{equation}
\sigma^{\%}_{Z} = \frac{\sigma_Z}{S_0}.
\end{equation}
This is a necessary but, as shown below, not sufficient correction.

\section{Error 3: Missing $\sqrt{\Delta t}$ scaling}
\label{sec:dtscale}

\texttt{HamiltonFilter}'s emission model treats its \texttt{regime\_volatilities} argument as the \emph{per-step return standard deviation} directly:
\begin{equation}
f(r_t \mid Z_t = k) = \N(r_t; 0, \sigma_k^{\text{step}}).
\end{equation}
But $\sigma^{\%}_Z$ from Section~\ref{sec:pctunits} is a \emph{volatility coefficient} in the SDE sense — it must still be scaled by $\sqrt{\Delta t}$ to produce the actual per-step standard deviation:
\begin{equation}
\sigma_k^{\text{step}} = \sigma^{\%}_k \sqrt{\Delta t}.
\end{equation}
With $\Delta t = 0.005$, $\sqrt{\Delta t} \approx 0.0707$; omitting this factor makes the assumed $\sigma_k^{\text{step}}$ about $14\times$ too large for \emph{both} regimes.

\paragraph{Why this is not a harmless uniform rescaling.} Although the \emph{ratio} $\sigma_1^{\text{step}}/\sigma_0^{\text{step}} = 10$ is preserved regardless of the missing factor, the correction is not scale-invariant for a fixed likelihood ratio evaluated at a fixed observed $r_t$: if $\sigma_0, \sigma_1$ are both inflated by the same multiplicative error, a genuine draw from regime 1's true (small) diffusion can end up numerically closer, in standardised units, to the (also-inflated) regime-0 Gaussian than to its own, destroying discrimination or even reversing it. Empirically, on a diffusion-only observation with $r_t \approx 3\times 10^{-5}$ genuinely drawn under regime 1:
\begin{itemize}
\item With the bug (missing $\sqrt{\Delta t}$): $f(r_t\mid Z{=}0) \approx 3810$, $f(r_t\mid Z{=}1) \approx 399$ — regime $0$ favoured by a factor of $\sim 9.6$, i.e.\ \textbf{backwards}.
\item Corrected: $f(r_t\mid Z{=}1) \approx 5167$, $f(r_t\mid Z{=}0)\approx 7.0$ — regime $1$ correctly favoured by a factor of $\sim 740$.
\end{itemize}
Classification accuracy on a representative 300-step episode rose from $0.363$ (worse than chance) to $0.900$ after this correction combined with Error 4 below.

This error was masked wherever the jump-mixture emission (Section~\ref{sec:jumpmix}) was also enabled, because the jump signal dominates the likelihood almost irrespective of the diffusion-only $\sigma$ scaling; it was fully exposed only in the two Gaussian-only filter configurations.

\section{Error 4: Jump-diffusion mixture emission and the limits of moment matching}
\label{sec:jumpmix}

\subsection{The mixture likelihood}

Conditional on regime $1$, a single step's return is itself a mixture over four elementary arrival outcomes (bid arrival $\times$ ask arrival $\in \{0,1\}^2$), each independently Bernoulli$(p)$ with $p$ from \eqref{eq:jumpprob}. Collapsing the two ``cancel'' outcomes (both or neither arrival) into a single zero-mean component gives
\begin{equation}
f(r_t \mid Z_t = 1) = w_0\, \N(r_t; 0, \sigma_1) + w_+\, \N(r_t; +\epsilon, \sigma_1) + w_-\, \N(r_t; -\epsilon, \sigma_1),
\label{eq:mixture}
\end{equation}
with
\begin{equation}
w_0 = (1-p)^2 + p^2, \qquad w_+ = w_- = p(1-p).
\end{equation}
In the symmetric-intensity case ($w_+ = w_-$), the mixture mean is zero, and by the law of total variance,
\begin{equation}
\Var(r_t \mid Z_t{=}1) = \E[\Var(r_t \mid \text{component})] + \Var(\E[r_t \mid \text{component}])
= \sigma_1^2 + (w_+ + w_-)\epsilon^2 = \sigma_1^2 + 2p(1-p)\epsilon^2.
\label{eq:mixturevar}
\end{equation}
This is exactly the variance used by the realised-variance (RV) feature's Gamma emission, $v_1 = \sigma_1^2 + 2p(1-p)\epsilon^2$.

\subsection{Why a moment-matched Gaussian is a poor substitute for the return likelihood}

One might ask whether it suffices to replace \eqref{eq:mixture} with a single Gaussian of matched variance,
\begin{equation}
f(r_t \mid Z_t{=}1) \overset{?}{\approx} \N\!\left(r_t; 0, \sqrt{\sigma_1^2 + 2p(1-p)\epsilon^2}\right).
\label{eq:momentmatch}
\end{equation}
This approximation is reasonable only when the separation ratio $\epsilon/\sigma_1$ is $O(1)$ or smaller, so that the three components blend into something unimodal. Here $\sigma_1 \approx 7.07\times 10^{-5}$ and $\epsilon = 0.01$, giving $\epsilon/\sigma_1 \approx 141$: the three components are essentially disjoint, narrow spikes, not a blurred continuum.

Quantitatively, with $\sigma_{\text{eff}} = \sqrt{\sigma_1^2 + 2p(1-p)\epsilon^2} \approx 6.5\times 10^{-3}$:
\begin{itemize}
\item At an actual jump return $r_t = \epsilon$: true mixture density $\approx 1185$; moment-matched Gaussian $\approx 18$ — the true model is $\sim 65\times$ more confident this is regime 1.
\item At $r_t \approx 0$ (no jump occurred): true mixture density $\approx 3272$; moment-matched Gaussian $\approx 62$ — the true ``no-jump'' state is far sharper than the smeared approximation suggests.
\end{itemize}
Matching only the first two moments discards almost all the information content of a single-step observation, because that information lives in \emph{which discrete outcome occurred}, not in a continuous spread around the mean. This is in contrast to the RV feature, which aggregates squared returns over an $8$--$10$-step window; there, the discreteness of individual jump/no-jump events is partially smoothed by aggregation (a mild central-limit effect), so the moment-matched variance \eqref{eq:mixturevar} is a defensible input to a Gamma approximation of the \emph{windowed} statistic even though it would not be an adequate replacement for the exact \emph{per-step} return likelihood \eqref{eq:mixture}.

\paragraph{Correction.} Two filter configurations (\texttt{plot\_final\_belief.py} and the belief/randomised agents in \texttt{simulate\_belief\_weighted.py}) were constructed with \texttt{jump\_size=None}, i.e.\ using \eqref{eq:momentmatch}-style plain-Gaussian emission with $v_1 = \sigma_1^2$ only (not even matching the true variance \eqref{eq:mixturevar}, since the jump contribution was not added at all in that code path). Enabling the true mixture \eqref{eq:mixture} (and thus also the corrected $v_1$ in the RV term) raised accuracy on the same representative episode from $0.363$ to $0.900$ (jointly with Error 3's correction).

\section{Error 5: Action ordering (bid/ask index convention)}
\label{sec:bidask}

\texttt{mbt\_gym}'s \texttt{LimitOrderModelDynamics} interprets the two-dimensional action vector as
\begin{equation}
a = (\delta^{b}, \delta^{a}), \qquad \texttt{BID\_INDEX}=0,\ \texttt{ASK\_INDEX}=1,
\end{equation}
confirmed directly from the fill/cash update rule (the sign convention on cash and inventory changes for each index matches the standard ``buy at bid, sell at ask'' convention only under this ordering).

The optimal-control code correctly computes $(\delta^{a,*}, \delta^{b,*})$ (verified numerically against \texttt{mbt\_gym}'s own reference \texttt{CarteaJaimungalMmAgent}, matching to four decimal places across a range of inventories), but every call site then constructed
\begin{equation}
a_{\text{sent}} = \big(\,\texttt{normalise}(\delta^{a,*}),\ \texttt{normalise}(\delta^{b,*})\,\big),
\end{equation}
i.e.\ the ask depth was sent into the bid slot and vice versa.

\paragraph{Consequence.} Since the naive baseline quotes symmetrically ($\delta^a = \delta^b$), the swap is a no-op for it — which is precisely why it was not the first place the bug was noticed. For the inventory-skewing agents (oracle, belief-weighted, randomised), the skew was applied to the \emph{wrong} side, actively working against inventory control rather than implementing it. Empirically, before correction the oracle agent's mean absolute inventory ($2.86$) was \emph{worse} than the naive baseline's ($2.32$); after correcting both this error and Error 6 (below), it improved to $1.68$ — clearly the best of the four agents, as theory would predict.

\section{Error 6: Terminal inventory aversion mismatch ($\alpha$)}
\label{sec:control}

The closed-form control is obtained via the standard Cartea--Jaimungal--Penalva linear-ODE method: define $w(t,q)$ solving
\begin{equation}
\frac{dw_q}{dt} = \kappa \phi q^2\, w_q - r_+\, w_{q-1} - r_-\, w_{q+1}, \qquad
w_q(T) = e^{-\kappa \alpha q^2},
\label{eq:hjb}
\end{equation}
with fill-intensity coefficients evaluated at the baseline (inventory-neutral, adverse-selection-adjusted) quote,
\begin{equation}
r_\pm = \lambda\, e^{-1 - \kappa \epsilon}.
\end{equation}
The optimal depths follow from finite differences of $\log w$:
\begin{equation}
\delta^{a,*}(t,q) = \frac1\kappa + \epsilon + \frac1\kappa \log\frac{w_{q+1}(t)}{w_q(t)}, \qquad
\delta^{b,*}(t,q) = \frac1\kappa + \epsilon + \frac1\kappa \log\frac{w_{q}(t)}{w_{q+1}(t)}.
\end{equation}

The environment's actual reward function, \texttt{RunningInventoryPenalty}, is exactly the discretisation of \eqref{eq:objective} with $\phi = \texttt{PER\_STEP\_INVENTORY\_AVERSION} = 0.01$ and $\alpha = \texttt{TERMINAL\_INVENTORY\_AVERSION} = 0$. The control code, however, hardcoded $\alpha = 0.001$ across every script — a value that does not correspond to any parameter of the actual environment being scored. Cross-checking \texttt{mbt\_gym}'s own reference agent confirms it reads $\alpha$ \emph{directly} from the reward function rather than hardcoding it, which is the correct pattern.

\paragraph{Consequence.} With $\alpha \ne 0$, \eqref{eq:hjb}'s terminal condition imposes an artificial pressure to liquidate inventory by $T$ that has no counterpart in the actual scored objective — the ``oracle'' was solving a subtly different control problem from the one it was being evaluated against.

\paragraph{Correction.} $\phi, \alpha$ are now imported directly from the environment's own constants rather than restated as separate literals.

\section{Error 7: Boundary-condition singularity in the inventory grid}
\label{sec:boundary}

The linear system \eqref{eq:hjb} is solved on a truncated inventory grid $q \in \{-Q_{\max}, \dots, Q_{\max}\}$, with a reflecting boundary condition implemented by explicitly zeroing the off-diagonal entries that would otherwise couple the boundary rows outward:
\begin{equation}
A_{0,1} = 0, \qquad A_{N-1,N-2} = 0.
\end{equation}
This decouples $w_{-Q_{\max}}(t)$ and $w_{Q_{\max}}(t)$ from the smooth interior recursion. Since the optimal depths are computed as log-ratios against the \emph{adjacent} grid row,
\begin{equation}
\delta^{a,*}(t,q) = \frac1\kappa + \epsilon + \frac1\kappa \log\frac{w_{q+1}(t)}{w_q(t)},
\end{equation}
exactly the two grid cells that reference a boundary row directly ($q = -Q_{\max}+1$ and $q = Q_{\max}$ for the ask depth; symmetrically for the bid depth) inherit that decoupling as a numerical singularity, while every other cell — including $q$ as close as $Q_{\max}-1$ — remains smooth and economically sensible (tight ask when long, wide ask when short, symmetric about $q=0$).

\paragraph{Empirical confirmation.} With $Q_{\max}=50$, $\kappa=2,\phi=0.01,\alpha=0$:
\begin{center}
\begin{tabular}{@{}rrr@{}}
\toprule
$q$ & $\delta^{a,*}(0,q)$ & \\
\midrule
$-48$ & $1.025$ & smooth \\
$-49$ & $65.294$ & \textbf{singular} (references boundary row) \\
\vdots & \vdots & \\
$49$ & $-0.025$ & smooth \\
$50$ & $-64.294$ & \textbf{singular} (references boundary row) \\
\bottomrule
\end{tabular}
\end{center}
Because realistic simulated inventory rarely approaches $Q_{\max}=50$ (mean $|q| \approx 1.7$--$2.3$ across all four agents), this artifact is invisible in ordinary simulation output but appears immediately in any diagnostic that reports $\min/\max$ depth over the full grid.

\paragraph{Correction.} Solve on a grid padded by one additional inventory level on each side ($Q_{\max}^{\text{solve}} = Q_{\max}+1$), then discard the outermost level before returning — pushing the singular cell outside the range ever consumed by \texttt{get\_control}, while leaving the public grid shape and range unchanged. Applied to all six independent copies of \texttt{build\_optimal\_control} in the codebase. This matters most for \texttt{optimal\_control.py} and \texttt{simulate\_oracle.py}, whose small $Q_{\max}=5$ grid puts the singular cell close enough to realistic inventory levels that it was plausibly reachable in practice, not just a display artifact — indeed, \texttt{optimal\_control.py}'s own sanity check (``min ask/bid depth should be $\ge 0$'') was silently failing before this fix.

\section{Why perfect regime information does not guarantee the oracle wins on raw mean PnL}
\label{sec:discussion}

After Errors 1--6 were corrected, the oracle agent achieved the lowest inventory and the best Sharpe ratio of the four agents, but \emph{not} the highest raw mean PnL relative to the belief-weighted agent. This is expected, for two distinct theoretical reasons.

\subsection{Risk-adjusted vs.\ raw objective}
The oracle is optimal for \eqref{eq:objective}, which explicitly trades expected spread capture against inventory variance via $\phi q_t^2$ and $\alpha q_T^2$. A metric that omits the penalty terms — such as $C_T - C_0$ (cash-only) — is scoring a \emph{different} functional than the one the control was derived to optimise. An agent that takes on more inventory risk can show a higher raw mean under such a metric while being strictly worse under \eqref{eq:objective}. Sharpe ratio (or, better, the realised value of \eqref{eq:objective} itself) is the metric under which the oracle's optimality claim is meaningful.

\subsection{The oracle is regime-static, not regime-switching-aware}
More subtly: \eqref{eq:hjb} assumes the regime $Z_t$ is fixed for the remainder of the episode. The environment's true regime process is a Markov chain that can switch \emph{within} an episode with non-trivial probability (given $P$'s off-diagonal entries). The oracle agent re-selects whichever regime-conditional solution matches the currently-true $Z_t$ at every step, but this is a myopic heuristic with respect to the \emph{true} switching dynamics — it never prices in the possibility that the regime may switch again before $T$.

A fully switching-aware control would instead solve a coupled system, e.g.\ for two regimes with switching intensities $\Lambda_{01}, \Lambda_{10}$,
\begin{equation}
\frac{dw_q^{(i)}}{dt} = \kappa \phi q^2\, w_q^{(i)} - r_+^{(i)} w_{q-1}^{(i)} - r_-^{(i)} w_{q+1}^{(i)} - \Lambda_{i,1-i}\left(w_q^{(1-i)} - w_q^{(i)}\right), \qquad i \in \{0,1\},
\label{eq:coupledhjb}
\end{equation}
with the two regimes' value functions linked by the transition-rate term. Solving \eqref{eq:coupledhjb} (rather than \eqref{eq:hjb} independently per regime) is a genuine extension beyond the scope of the corrections in this note, and is flagged here as a modelling limitation rather than a bug: the current oracle has \emph{perfect state information} but does not use it to solve the \emph{fully specified} control problem implied by the environment's actual regime dynamics. Blackwell-type information-dominance results (more information cannot reduce attainable expected utility) presuppose optimal use of that information for the objective being scored; a regime-static approximation, however accurate within a single regime spell, does not inherit that guarantee once mid-episode switching is a first-order effect.

\paragraph{Correction (post hoc).} \eqref{eq:coupledhjb} above was written speculatively, before the coupled system was ever actually solved, and its coupling term $\Lambda_{i,1-i}(w_q^{(1-i)}-w_q^{(i)})$ is \emph{not} correct: it assumes the switching contribution linearises under the same $w=e^{\kappa h}$ transformation used for the single-regime HJB, which is not justified and (as verified in Section~\ref{sec:switchingaware}) is not the case. The switching term is additive and linear in $h$ (via the value-function ansatz $V^{(i)}=x+qs+h_q^{(i)}(t)$, under which the switching contribution is exactly $\nu_{i,1-i}(h_q^{(1-i)}-h_q^{(i)})$, since cash/inventory/price do not jump when the regime switches), but this does \emph{not} carry over to a simple linear term in $w$ once re-expressed via $w=e^{\kappa h}$ -- the correct $w$-space equation would carry a $\log$-ratio inside the coupling term, not a plain difference of $w$'s. Section~\ref{sec:switchingaware} solves the correct coupled system directly in $h$-space and reports the resulting quantitative answer to the question raised in this subsection.

\subsection{The randomised agent}
Separately, the randomised control (Bernoulli-sampling a hard regime label from the belief each step, cf.\ Section 3.2.3 of the main text) underperforms the belief-weighted interpolation (Section 3.2.2) even after all corrections. This is an expected property of the design, not a bug: collapsing a continuous belief into a discrete draw discards information relative to smooth interpolation of the two regime-conditional controls, adding sampling noise without a compensating benefit.

\section{Extension: stochastic (Exponential) jump size}
\label{sec:expjump}

\subsection{Motivation}
Under the original specification, \eqref{eq:midprice}'s jump magnitude $\epsilon$ was a fixed constant applied on every market-order arrival. With the calibration used throughout this note ($\kappa=1.5$, $\epsilon=0.5$, so $1/\kappa \approx 0.667 > \epsilon$), this has a structural consequence: since every agent's quoted depth satisfies $\delta \ge 1/\kappa$ by construction, and the single-trade markout on a filled quote is $\delta - \epsilon$, \emph{no individual fill can ever lose money to adverse selection} -- the baseline inventory-neutral half-spread alone already exceeds the entire deterministic jump. Empirically, this produced policy-comparison runs in which all four agents' episode PnL was reliably positive (loss frequency $<2\%$ across 1000-episode samples), which is not a realistic picture of adverse-selection risk.

\subsection{Change}
The jump magnitude on each arrival is now random: $\mathrm{Exponential}(\text{mean}=\epsilon)$, independently drawn per arrival, rather than the fixed constant $\epsilon$. This only requires reinterpreting $\epsilon$ as $\E[\text{jump magnitude}]$ throughout -- which \eqref{eq:hjb}'s coefficient $r_\pm = \lambda e^{-1-\kappa\epsilon}$ already assumes, since it is derived from evaluating the fill intensity at the \emph{expected} optimal quote. Consequently \textbf{no change to \texttt{build\_optimal\_control} (Errors 1--7) or the closed-form control was required} -- the HJB solution is unaffected by this randomisation to leading order (a small Jensen-gap correction, $\E[e^{-\kappa\epsilon}]\ge e^{-\kappa\E[\epsilon]}$, is second-order and not corrected here).

What \emph{does} need to change is the Hamilton filter's regime-1 return emission (Section~\ref{sec:jumpmix}), since it must now describe the likelihood of returns generated by a random-magnitude jump, not a fixed one.

\subsection{Corrected emission derivation}
Conditional on regime 1, each step has four mutually exclusive arrival outcomes: neither side arrives (prob $(1-p)^2$, no jump); one side arrives (prob $p(1-p)$ each), contributing a one-sided jump $A\sim\mathrm{Exponential}(\epsilon)$; or both sides arrive (prob $p^2$), contributing a net jump $A-B$ for independent $A,B\sim\mathrm{Exponential}(\epsilon)$.

Unlike the fixed-jump case, $A-B$ does \emph{not} collapse to zero: the difference of two i.i.d.\ $\mathrm{Exponential}(\epsilon)$ random variables is exactly $\mathrm{Laplace}(0,\epsilon)$, which is itself a $50/50$ mixture of $+\mathrm{Exponential}(\epsilon)$ and $-\mathrm{Exponential}(\epsilon)$. So the "both arrive" mass splits evenly into the same two branches as the one-sided case, giving combined weight $p(1-p) + p^2/2 = p(1-p/2)$ per branch, and
\begin{equation}
f(r_t \mid Z_t{=}1) = (1-p)^2\,\N(r_t;0,\sigma_1) \;+\; p\Big(1-\tfrac{p}{2}\Big)\Big[\mathrm{EMG}(r_t;\sigma_1,\epsilon) + \mathrm{EMG}(-r_t;\sigma_1,\epsilon)\Big],
\label{eq:emgmixture}
\end{equation}
where $\mathrm{EMG}(\cdot;\sigma,\epsilon)$ is the density of $\mathrm{Normal}(0,\sigma) + \mathrm{Exponential}(\epsilon)$ (an Exponentially Modified Gaussian, \texttt{scipy.stats.exponnorm} with shape $K=\epsilon/\sigma$). The three weights sum to 1 by construction. This replaces the fixed-jump mixture of Section~\ref{sec:jumpmix}, \eqref{eq:mixture}. The realised-variance emission's target variance is correspondingly corrected from $v_1 = \sigma_1^2 + 2p(1-p)\epsilon^2$ to
\begin{equation}
v_1 = \sigma_1^2 + 2p(2-p)\epsilon^2,
\label{eq:v1exp}
\end{equation}
reflecting the extra variance contributed by the jump \emph{magnitude} itself (each $\mathrm{Exponential}(\epsilon)$ draw has $\E[A^2]=2\epsilon^2$, vs.\ a deterministic jump's $\epsilon^2$).

\subsection{Validation}
Three checks confirm the implementation matches the derivation:
\begin{enumerate}
\item \textbf{Jump mechanics.} Isolating one-sided arrivals in \texttt{ArrivalJumpMidpriceModel} against $\mathrm{Exponential}(\text{mean}=\epsilon)$: KS $p=0.90$. Isolating both-sided arrivals against $\mathrm{Laplace}(0,\epsilon)$: KS $p=0.94$ ($n=200{,}000$ each).
\item \textbf{Filter emission vs.\ simulated returns.} 500,000 i.i.d.\ regime-1 returns simulated from the actual midprice model (realistic $\mathrm{Bernoulli}(p)$ arrivals) tested against \eqref{eq:emgmixture}'s exact closed-form CDF: KS $\mathrm{stat}=0.00134$, $p=0.656$, consistent across subsample sizes. (An earlier attempt using a numerically-integrated CDF on a coarse grid spuriously rejected at $p\approx0$ -- an artefact of integration error, not misspecification; resolved by using the exact CDF, $\N$'s and \texttt{exponnorm}'s CDFs in closed form.)
\item \textbf{End-to-end filter accuracy} on the same representative episode used throughout this note improved from 0.83 to 0.97 after the correction -- the old fixed-jump emission was misspecified once the true generative process became random.
\end{enumerate}

\subsection{Empirical effect on policy comparison}
Loss frequency (fraction of 1000 episodes with negative markout PnL) and worst-case episode PnL, before vs.\ after:
\begin{center}
\begin{tabular}{@{}lrrrr@{}}
\toprule
Agent & Loss freq.\ (fixed $\epsilon$) & Loss freq.\ (Exp.\ $\epsilon$) & Min PnL (fixed) & Min PnL (Exp.) \\
\midrule
Oracle & 1.5\% & 5.4\% & $-15.10$ & $-109.69$ \\
Belief-weighted & 1.0\% & 6.0\% & $-39.79$ & $-60.24$ \\
Randomised & 0.9\% & 5.4\% & $-21.52$ & $-96.79$ \\
Naive & 2.1\% & 8.3\% & $-48.09$ & $-151.55$ \\
\bottomrule
\end{tabular}
\end{center}
Mean PnL is essentially unchanged (the Exponential's mean is still $\epsilon$), but the loss tail is $3$--$4\times$ more frequent and $2$--$7\times$ deeper, and naive -- which does not skew quotes with inventory or belief -- now clearly has both the highest loss frequency and the largest realised losses of the four agents, consistent with theory. Sharpe ratios on a smaller 500-episode check also now separate cleanly (Oracle $0.95>$ Randomised $0.89>$ Belief-weighted $0.87>$ Naive $0.84$), rather than clustering tightly together as under the fixed-jump model.

\section{Summary of empirical validation}

\begin{table}[h]
\centering
\begin{tabular}{@{}lrr@{}}
\toprule
Metric & Before & After \\
\midrule
\texttt{plot\_final\_belief.py} accuracy & 0.363 & 0.900 \\
\texttt{plot\_filter\_comparison.py} accuracy & --- & 0.891 \\
\texttt{plot\_spread\_diagnostic.py} accuracy & --- & 0.844 \\
\texttt{test\_hamilton\_filter.py} belief separation & 0.075 & 0.406--0.530 \\
Oracle mean $|q|$ (inventory) & 2.86 (worse than naive) & 1.67--1.73 (best of four) \\
Belief-weighted mean $|q|$ & 2.58 & 1.66 \\
Randomised mean $|q|$ & 2.51 & 1.66 \\
Naive mean $|q|$ & 2.29--2.32 & 2.31 \\
Oracle Sharpe & --- & best of four (0.136) \\
\bottomrule
\end{tabular}
\caption{Before/after comparison across the corrected pipeline. ``Before'' entries are omitted where the pre-fix configuration was not separately measured under an equivalent metric.}
\end{table}

\section{Files touched}

\begin{itemize}
\item \texttt{envs/regime\_env.py} --- added \texttt{raw\_state}/\texttt{raw\_cash}/\texttt{raw\_inventory}/\texttt{raw\_midprice} accessors and \texttt{info['raw\_state']} (Error 1).
\item \texttt{envs/make\_envs.py} --- added \texttt{R0\_VOLATILITY\_PCT}, \texttt{R1\_VOLATILITY\_PCT}, \texttt{EPSILON\_PCT} (Error 2) and \texttt{R0\_SIGMA\_STEP}, \texttt{R1\_SIGMA\_STEP} (Error 3).
\item \texttt{beliefs/hamilton\_filter.py} --- documentation clarification only; emission logic unchanged (the bugs were in how it was \emph{called}, not in the class itself).
\item \texttt{test\_hamilton\_filter.py}, \texttt{simulate\_belief\_weighted.py}, \texttt{simulate\_oracle.py}, \texttt{plot\_belief\_evolution.py}, \texttt{plot\_filter\_comparison.py}, \texttt{plot\_spread\_diagnostic.py}, \texttt{plot\_final\_belief.py} --- updated to use the corrected accessors/constants, the corrected action ordering (Error 5), and the environment-matched $\phi,\alpha$ (Error 6); \texttt{plot\_final\_belief.py} and \texttt{simulate\_belief\_weighted.py}'s \texttt{make\_filter()} additionally switched from Gaussian-only to jump-mixture emission (Error 4).
\item \texttt{simulate\_belief\_weighted.py}, \texttt{simulate\_oracle.py}, \texttt{plot\_final\_belief.py}, \texttt{plot\_filter\_comparison.py}, \texttt{plot\_spread\_diagnostic.py}, \texttt{optimal\_control.py} --- each independent copy of \texttt{build\_optimal\_control} updated with the padded-grid boundary fix (Error 7).
\item \texttt{envs/arrival\_jump\_midprice.py} --- jump magnitude on each arrival changed from the fixed constant \texttt{jump\_size} to an independent $\mathrm{Exponential}(\text{mean}=\texttt{jump\_size})$ draw (Section~\ref{sec:expjump}).
\item \texttt{beliefs/hamilton\_filter.py} --- regime-1 return emission replaced with the EMG-based mixture \eqref{eq:emgmixture}; realised-variance target \eqref{eq:v1exp} corrected accordingly (Section~\ref{sec:expjump}).
\item \texttt{simulate\_belief\_weighted.py} --- added mark-to-market PnL (previously cash-only, silently excluding terminal inventory value), a risk-adjusted objective column (cumulative environment reward, matching \eqref{eq:objective}), and \%-of-notional-capital normalisation for all reported PnL/objective figures.
\item \texttt{envs/make\_envs.py} --- \texttt{make\_regime\_envs} gained optional \texttt{epsilon}, \texttt{r0\_volatility}, \texttt{r1\_volatility} keyword arguments (each defaulting to the existing module-level constant, so all existing call sites are unaffected), enabling the epsilon sweeps of Sections~\ref{sec:epsilonsweep}--\ref{sec:policyepsilonsweep} without duplicating the environment-construction logic.
\end{itemize}

\paragraph{Current environment calibration.} As of this note's latest revision: $\kappa=1.5$ (\texttt{KAPPA}, calibrated to the Avellaneda--Stoikov (2008) fill-intensity specification), $\epsilon=0.5$ (\texttt{EPSILON}, now the \emph{mean} jump size per Section~\ref{sec:expjump}), $\lambda=140$ (\texttt{LAMBDA}), $\sigma_0^{\text{pct}}=\sigma_1^{\text{pct}}=0.01$ (\texttt{R0\_VOLATILITY}/\texttt{R1\_VOLATILITY}, equal since the recalibration of Section~\ref{sec:equalvolrecal} -- previously $\sigma_1^{\text{pct}}=0.03$, a $3\times$ ratio), $\phi=0.01$, $\alpha=0.001$ (\texttt{PER\_STEP\_INVENTORY\_AVERSION}/\texttt{TERMINAL\_INVENTORY\_AVERSION}), all in \texttt{envs/make\_envs.py}. Note $\kappa$ differs from the $\kappa=2$ used in the worked examples of Sections~\ref{sec:pctunits}--\ref{sec:boundary} above, which describe the state of the environment at the time those specific bugs were diagnosed, prior to recalibrating $\kappa$ against Avellaneda--Stoikov.

\appendix

\section{Comprehensive specification of the current environment (for cross-checking)}
\label{sec:currentstate}

This appendix is a self-contained, current-state description of exactly how the simulation runs today, independent of the historical bug-fix narrative in Sections~\ref{sec:norm}--\ref{sec:expjump} above. It is intended as a single reference against which to cross-check the code for further errors.

\subsection{Architecture}
\texttt{RegimeSwitchingEnv} (\texttt{envs/regime\_env.py}) wraps two independent \texttt{mbt\_gym} \texttt{TradingEnvironment} instances, one per regime, each with its own midprice model, arrival model, fill-probability model and reward function. At any time exactly one sub-environment is ``active''; \texttt{RegimeSwitchingEnv.step()} delegates to the active sub-environment, then (if \texttt{switch\_within\_episode=True}, the configuration used throughout this note) draws the next regime from the Markov chain and, on an actual switch, copies cash/inventory/price into the newly-active sub-environment's state (\texttt{\_sync\_state}) so the two regimes share one continuous portfolio and price path rather than each resuming from its own frozen initial state. Each sub-environment keeps its own internal clock (time is not synced), so \texttt{done} fires when either sub-environment reaches its own \texttt{n\_steps}.

The true regime $Z_t$ is exposed via \texttt{info['true\_regime']} and the \texttt{current\_regime} property (for the oracle and for scoring/diagnostics only); it is never fed to the belief-driven agents, which instead observe only the midprice series via a \texttt{HamiltonFilter}.

\subsection{Regime process}
$Z_t \in \{0,1\}$ evolves as a discrete-time Markov chain with row-stochastic transition matrix
\begin{equation}
P = \begin{pmatrix} 0.95 & 0.05 \\ 0.02 & 0.98 \end{pmatrix},
\label{eq:sys-P}
\end{equation}
whose stationary distribution is $\pi^{\text{stat}} = (0.286, 0.714)$ -- i.e.\ the environment spends roughly $71\%$ of stationary time in the adverse-selection regime, not $50\%$.

\subsection{Midprice dynamics}
Regime 0 (\texttt{BrownianMotionMidpriceModel}): pure diffusion, $S_{t+\Delta t} = S_t + \sigma_0 \sqrt{\Delta t}\,\xi_t$.

Regime 1 (\texttt{ArrivalJumpMidpriceModel}, \texttt{envs/arrival\_jump\_midprice.py}):
\begin{equation}
S_{t+\Delta t} = S_t + \sigma_1\sqrt{\Delta t}\,\xi_t + A_t\,\mathbb{1}[\text{buy MO arrival}] - B_t\,\mathbb{1}[\text{sell MO arrival}],
\label{eq:sys-midprice1}
\end{equation}
where $A_t, B_t \overset{\text{iid}}{\sim} \mathrm{Exponential}(\text{mean}=\epsilon)$ are drawn fresh every step (regardless of whether that side actually arrives, mirroring the unconditional diffusion draw), and $\xi_t\sim\N(0,1)$ independently. This is the corrected, random-jump-size version of \eqref{eq:midprice}; see Section~\ref{sec:expjump} for the derivation and validation. The jump fires on \emph{arrivals}, not on the market maker's own fills -- i.e.\ the price moves whenever informed order flow shows up, whether or not that particular order happened to execute against this agent's quote.

\subsection{Arrival and fill mechanics -- two independent processes}
\label{sec:sys-arrivalfill}
Per side (bid/ask), each step draws two \emph{separate}, independently-seeded Bernoulli random variables:
\begin{align}
\text{arrival} &\sim \mathrm{Bernoulli}(p), \quad p = \lambda\Delta t, &&\text{(\texttt{PoissonArrivalModel.get\_arrivals}, own RNG)} \label{eq:sys-arrival}\\
\text{fill} &\sim \mathrm{Bernoulli}\big(e^{-\kappa\delta}\big), &&\text{(\texttt{ExponentialFillFunction.get\_fills}, separate RNG)} \label{eq:sys-fill}
\end{align}
where $\delta$ is the agent's quoted depth on that side. Critically, \eqref{eq:sys-fill} does not condition on \eqref{eq:sys-arrival} at the point of drawing -- it is a fresh uniform draw compared against a depth-dependent probability, independent of whether an order actually arrived. \texttt{make\_envs.py} seeds the two models from independent child seeds (\texttt{SeedSequence.spawn}), so they are genuinely statistically independent streams, not merely computed separately from a shared source.

An actual executed trade requires \emph{both}: \texttt{ModelDynamics.update\_state} applies inventory/cash changes scaled by \texttt{arrivals * fills} (elementwise product), i.e.\ a side only trades when a market order arrived \emph{and} the independent fill coin-flip succeeded. Quoting a wider $\delta$ only reduces the probability of \emph{being} the counterparty (via \eqref{eq:sys-fill}); it has no effect on whether the arrival happens or on the midprice jump in \eqref{eq:sys-midprice1}, which fires on arrivals alone. Adverse selection therefore specifically afflicts the subset of steps where both an arrival \emph{and} a fill occur on the same side, not all arrivals.

\subsection{Reward, markout PnL, and the risk-adjusted objective}
The reward function is \texttt{RunningInventoryPenalty} (\texttt{mbt\_gym.rewards.RewardFunctions}), which computes at each step the mark-to-market value change $d(C_t + q_tS_t)$ minus the running penalty $\phi q_t^2\Delta t$, minus $\alpha q_T^2$ only on the terminal step -- i.e.\ exactly a discretisation of \eqref{eq:objective}. Two distinct quantities are reported per episode in \texttt{simulate\_belief\_weighted.py}:
\begin{align}
\mathrm{PnL}_T &= \big(C_T + q_TS_T - \alpha q_T^2\big) - \big(C_0+q_0S_0\big) &&\text{(markout PnL, excludes running penalty)}\label{eq:sys-pnl}\\
\mathrm{Obj}_T &= \sum\nolimits_t \text{reward}_t \;=\; \mathrm{PnL}_T - \phi\!\sum\nolimits_t q_t^2\Delta t &&\text{(cumulative environment reward} = J\text{ sample)}\label{eq:sys-obj}
\end{align}
Both are reported as a percentage of $K = S_0 \times q_{\max}$ ($q_{\max}=10{,}000$, \texttt{TradingEnvironment}'s hard inventory ceiling), i.e.\ the dollar value of the largest position the agent could ever hold -- \emph{not} a deterministic constraint any agent approaches in practice (realised mean $|q|\approx2$--$3$), just a fixed capital-at-risk denominator so PnL/objective are comparable across configurations. Sharpe ratios are computed on the raw (pre-\%) values and are identical either way (scale-invariant).

\subsection{Hamilton filter}
\texttt{beliefs/hamilton\_filter.py} tracks $\pi_t = P(Z_t=1\mid\mathcal{F}_t)$ via a Hamilton (1989) recursion with $r$ lagged states in the joint posterior ($r=1$ by default), operating on the percentage return \eqref{eq:pctreturn}. Regime 0's emission is $\N(0,\sigma_0^{\text{step}})$; regime 1's is the EMG-based mixture \eqref{eq:emgmixture} derived in Section~\ref{sec:expjump}, matched to the random-jump-size midprice model of \eqref{eq:sys-midprice1}. An optional rolling realised-variance emission augmentation exists in the code (\texttt{rv\_window}, \texttt{rv\_weight}) but is \textbf{disabled by default} -- it double-counts each return's evidence across up to \texttt{rv\_window} consecutive steps (the same observation contributes to its own step's return likelihood \emph{and} to the next \texttt{rv\_window}$-1$ steps' RV likelihood), which was found to collapse belief fluctuation within a true regime-1 spell (min belief $1.0$ throughout, vs.\ a legitimate dip to $0.24$ without it). It is retained only for ablation scripts (\texttt{plot\_belief\_evolution.py}); do not enable it for anything feeding a trading decision or a reported accuracy figure.

\subsection{Optimal control and the four policies}
The closed-form control \eqref{eq:hjb} is solved independently per regime (regime-static, not regime-switching-aware -- Section~\ref{sec:discussion}), on an inventory grid padded by one level per side to avoid the boundary singularity of Section~\ref{sec:boundary}. Four policies are compared, all acting through the same $(\delta^{a,*},\delta^{b,*})$ tables:
\begin{itemize}
\item \textbf{Oracle}: uses \texttt{env.current\_regime} directly (perfect, instantaneous regime knowledge) to select which regime's control table to query.
\item \textbf{Belief-weighted}: queries \emph{both} regimes' tables at the current $(t,q)$ and linearly interpolates by the filter's belief, $\delta_t = (1-\pi_t)\delta^{(0)}_t + \pi_t\delta^{(1)}_t$.
\item \textbf{Randomised}: samples a hard regime label $I_t\sim\mathrm{Bernoulli}(\pi_t)$ each step and fully commits to that regime's control.
\item \textbf{Naive}: quotes a fixed, non-inventory-skewing depth $\delta = 1/\kappa + \epsilon/2$ on both sides every step, regardless of regime, belief, or inventory.
\end{itemize}
(An ad hoc fifth variant -- always quoting regime 0's optimal control regardless of true regime -- was tested in this session as a diagnostic. It matched oracle's inventory-skewing behaviour, since that term depends only on $\kappa,\phi,\alpha$ -- shared across regimes -- but lost the $\epsilon$ adverse-selection padding during the $71\%$ of time spent in regime 1, giving it a higher loss frequency than naive [$7.8\%$ vs.\ $6.4\%$ per $500$ episodes] despite a similar or better Sharpe ratio. Not part of the main four-way comparison.)

\subsection{Current full parameter table}
\label{sec:paramtable}
\begin{center}
\begin{tabular}{@{}lll@{}}
\toprule
Symbol & Value & Code constant (\texttt{envs/make\_envs.py} unless noted) \\
\midrule
$T$ (terminal time) & $1.0$ & \texttt{TERMINAL\_TIME} \\
$N$ (steps) & $4{,}000$ & \texttt{N\_STEPS} (refined from $200$; Section~\ref{sec:discretisation}) \\
$\Delta t$ & $0.00025$ & \texttt{STEP\_SIZE} (refined from $0.005$) \\
$\kappa$ (fill decay) & $1.5$ & \texttt{KAPPA} (Avellaneda--Stoikov 2008 calibration) \\
$\sigma_0^{\text{pct}}$ & $0.01$ & \texttt{R0\_VOLATILITY} (absolute; $/100$ for pct.) \\
$\sigma_1^{\text{pct}}$ & $0.01$ (equal to $\sigma_0$; was $0.03$ -- Section~\ref{sec:equalvolrecal}) & \texttt{R1\_VOLATILITY} \\
$\epsilon$ (mean jump size) & $0.5$ & \texttt{EPSILON} (absolute price units) \\
$\lambda$ (arrival intensity, both sides) & $140$ & \texttt{LAMBDA} (unchanged by the discretisation refinement) \\
$\phi$ (running inventory aversion) & $0.01$ & \texttt{PER\_STEP\_INVENTORY\_AVERSION} \\
$\alpha$ (terminal inventory aversion) & $0.001$ & \texttt{TERMINAL\_INVENTORY\_AVERSION} \\
$P$ (transition matrix, $\Delta t=0.00025$) & see \eqref{eq:Pnew} & \texttt{TRANSITION\_MATRIX} (rescaled from \eqref{eq:sys-P}; Section~\ref{sec:discretisation}) \\
$S_0$ (initial price) & $100.0$ & \texttt{INITIAL\_PRICE} \\
$q_{\max}$ (hard inventory ceiling) & $10{,}000$ & \texttt{TradingEnvironment} default \\
$Q_{\max}$ (control-grid truncation) & $50$ & per-script \texttt{Q\_MAX} (not a real constraint; see Section~\ref{sec:boundary}) \\
$r$ (filter lagged states) & $1$ & \texttt{HamiltonFilter} default \\
RV emission & disabled & \texttt{rv\_window=0, rv\_weight=0.0} default \\
\bottomrule
\end{tabular}
\end{center}
Note: $P$ in \eqref{eq:sys-P} above and the transition-duration discussion in Section~\ref{sec:sys-arrivalfill}'s parent subsections describe the \emph{original} $\Delta t=0.005$ calibration; the live environment now uses the rescaled $\Delta t=0.00025$ matrix \eqref{eq:Pnew} derived in Section~\ref{sec:discretisation}, which preserves the same continuous-time generator and hence the same stationary distribution and (to a much closer approximation) the same expected regime durations.

\subsection{Known open items for cross-checking}
\begin{itemize}
\item The Jensen's-gap approximation flagged in Section~\ref{sec:expjump}: \eqref{eq:hjb} uses $e^{-\kappa\E[\epsilon]}$ where the fully correct term under a random jump is $\E[e^{-\kappa\epsilon}] \ge e^{-\kappa\E[\epsilon]}$. Not corrected; likely small given $\kappa\sigma_1 \ll 1$, but not quantified.
\item The oracle solves the regime-\emph{static} HJB (Section~\ref{sec:discussion}). \textbf{Resolved as a sensitivity analysis, not adopted into production}: a fully switching-aware coupled HJB \eqref{eq:switchinghjb} was implemented and solved (Section~\ref{sec:switchingaware}), and shown \emph{not} to explain the oracle-naive similarity -- the switching-aware and static controls differ only at inventory levels the simulation rarely visits.
\item Five independent copies of a callable \texttt{build\_optimal\_control(kappa, lam, eps, phi, alpha, ...)} function exist (\texttt{simulate\_belief\_weighted.py}, \texttt{simulate\_oracle.py}, \texttt{plot\_final\_belief.py}, \texttt{plot\_spread\_diagnostic.py}, \texttt{plot\_filter\_comparison.py}); a fix (or further bug) applied to one is not automatically applied to the others. \texttt{optimal\_control.py} is \emph{not} a sixth copy of the function -- it inlines the same derivation at module scope (not as a callable), using its own hardcoded example parameters ($\kappa=2.0,\lambda=0.75,\epsilon=0.3,\phi=10^{-5}$) that no longer match the live environment calibration in the table above ($\kappa=1.5,\lambda=140,\epsilon=0.5,\phi=0.01$); it independently includes the same padded-grid boundary fix, but should not be read as reflecting current parameters.
\item \texttt{fracNeg} (episode loss frequency) and tail severity are sensitive to the specific jump-size distribution family chosen (Exponential here); this was a deliberate modelling choice, not derived from data, and other heavy-tailed choices (e.g.\ log-normal, Pareto) would give different loss-tail behaviour for the same mean $\epsilon$.
\item The realised-variance filter emission augmentation remains implemented but disabled; if re-enabled for any future ablation, its known double-counting failure mode should be re-verified against the corrected EMG return emission (it predates that correction).
\end{itemize}

\section{Time-discretisation refinement: $\Delta t = 0.005 \to 0.00025$}
\label{sec:discretisation}

\subsection{Motivation}
The original discretisation used $\lambda\Delta t = 140 \times 0.005 = 0.7$ as the per-step, per-side arrival probability. This is an $O(1)$ quantity, not small, so the standard continuous-time Poisson approximation underlying the HJB/DPP derivation -- $P(\text{exactly one arrival in }dt) = \lambda\, dt + o(dt)$, $P(\text{two-or-more arrivals in }dt) = o(dt)$ -- was badly violated: empirically, $P(\text{both sides arrive in the same step}) = 0.7^2 = 0.49$ was the \emph{same order of magnitude} as $P(\text{exactly one arrival}) = 2(0.7)(0.3) = 0.42$, not two orders smaller as the continuous-time idealisation assumes.

\subsection{Change}
\begin{equation}
T = 1.0, \qquad n_{\text{steps}} = 4{,}000, \qquad \Delta t = T/n_{\text{steps}} = 0.00025, \qquad \lambda = 140 \ \text{(unchanged)},
\end{equation}
giving $\lambda\Delta t = 0.035$, now genuinely $\ll 1$. $\lambda$ itself is not altered -- only the step size.

\subsection{Transition-matrix rescaling}
Simply reusing the original $\Delta t_{\text{old}}=0.005$ matrix \eqref{eq:sys-P} at the new $\Delta t$ would make regime switching $20\times$ faster in continuous-time terms (the same per-step probabilities now apply to $20\times$ as many steps per unit time). Instead, the continuous-time generator implied by \eqref{eq:sys-P} is recovered and re-discretised at the new $\Delta t$:
\begin{equation}
Q = \frac{\log P_{\text{old}}}{\Delta t_{\text{old}}}, \qquad P_{\text{new}} = \exp(Q\, \Delta t_{\text{new}}),
\label{eq:Pnew}
\end{equation}
using the matrix logarithm/exponential. Concretely,
\begin{equation}
Q = \begin{pmatrix} -10.367242 & 10.367242 \\ 4.146897 & -4.146897 \end{pmatrix}, \qquad
P_{\text{new}} = \begin{pmatrix} 0.997413 & 0.002587 \\ 0.001035 & 0.998965 \end{pmatrix}.
\end{equation}
$Q$ has no imaginary component (verified: $P_{\text{old}}$ has two real, positive eigenvalues, $1$ and $0.93$), and $P_{\text{new}}$'s rows already sum to $1$ to floating-point precision (renormalised defensively regardless). Writing $\nu_{01}=Q_{01}=10.367242$, $\nu_{10}=Q_{10}=4.146897$ for the continuous-time transition intensities, the implied mean regime durations are $1/(-Q_{00})=0.096458$ (regime $0$) and $1/(-Q_{11})=0.241144$ (regime $1$), and the stationary distribution $\pi = (0.2857, 0.7143)$ is preserved \emph{exactly} by construction (both $P_{\text{old}}$ and $P_{\text{new}}$ share the same generator $Q$, hence the same null space).

\subsection{Grid and indexing corrections}
Two related but distinct fixes accompanied the step-size change:
\begin{itemize}
\item \textbf{Control-table time grid.} All five copies of \texttt{build\_optimal\_control} previously built their time grid via \texttt{np.linspace(0, T, n\_t)}, which spaces points by $T/(n_t-1)$, not $\Delta t = T/n_t$ -- a subtle off-by-one mismatch against the actual simulation step size (small at $n_t=200$, but worth eliminating rigorously at $n_t=4{,}000$). Changed to \texttt{np.arange(n\_t) * dt}, giving \texttt{t\_grid[n] = n*dt} exactly.
\item \textbf{Global time index.} The same five scripts' main loops previously reconstructed the oracle's control-table time index from the \emph{normalised} time observation (\texttt{time\_frac = (obs\_flat[2]+1)/2; t\_idx = round(time\_frac*(N\_T-1))}). Replaced with \texttt{t\_idx = env.current\_step} -- the wrapper's authoritative global step counter (Section~\ref{sec:currentstate}'s clock refactor), avoiding an unnecessary floating-point round-trip through the normalised observation.
\end{itemize}
Diffusion scaling ($\sigma\sqrt{\Delta t}$, via \texttt{R0\_SIGMA\_STEP}/\texttt{R1\_SIGMA\_STEP}) and the running-inventory-penalty's $\Delta t$ multiplier both propagate automatically from the new \texttt{STEP\_SIZE} constant -- verified directly (ratio of new to old $\sigma$ exactly $\sqrt{1/20}$; the reward function's internally-computed $\Delta t$, read from consecutive states' time stamps, equals $0.00025$ on every sampled step) rather than assumed.

\subsection{Empirical validation}
\paragraph{Arrival-count distribution} (600,000 steps, naive policy, both regimes pooled):
\begin{center}
\begin{tabular}{@{}lrr@{}}
\toprule
Quantity & Empirical & Theoretical \\
\midrule
$P(\text{buy arrival})$ & $0.03517$ & $0.035$ \\
$P(\text{sell arrival})$ & $0.034728$ & $0.035$ \\
$P(\text{both arrive})$ & $0.001202$ & $0.0012250$ \\
$P(\text{exactly one})$ & $0.067495$ & $0.06755$ \\
$\E[\text{buy arrivals per episode}]$ & $140.68$ & $140$ \\
$\E[\text{sell arrivals per episode}]$ & $138.91$ & $140$ \\
\bottomrule
\end{tabular}
\end{center}
All in close agreement -- $P(\text{both arrive})$ has fallen from $0.49$ (old $\Delta t$) to $0.0012$, a genuinely second-order event as the theory requires. A dedicated figure (\texttt{arrival\_discretisation\_verification\_*.png}) compares old vs.\ new empirical/theoretical bars for all three categories side by side.

\paragraph{Regime-duration re-estimation.} A first attempt estimated mean dwell time by naively averaging all contiguous same-regime segment lengths within fixed $4{,}000$-step episodes. This is biased: segments truncated by the episode boundary are under-counted, but \emph{excluding} boundary segments entirely makes the bias \emph{worse}, not better, because short episodes (by chance) fit more complete cycles and so contribute disproportionately more ``interior'' segments than long ones -- a length-biased-sampling artefact, not a code error (confirmed by reproducing the same pattern in a bare Markov-chain simulation using only \eqref{eq:Pnew}, with no environment involved). The statistically correct estimator for censored dwell times is the standard survival-analysis MLE,
\begin{equation}
\widehat{\mu}_i = \frac{\text{total time observed in state } i}{\text{number of \emph{observed} exits from state } i},
\end{equation}
which uses all occupied time (no bias) while only counting actually-realised transitions in the denominator. Applied to $150$ episodes ($n_{\text{exits},0}=466$, $n_{\text{exits},1}=462$):
\begin{center}
\begin{tabular}{@{}lrrr@{}}
\toprule
& MLE (empirical) & Continuous-time target & New-$\Delta t$ discrete-implied \\
\midrule
Regime 0 mean duration & $0.091758$ & $0.096458$ & $0.096633$ \\
Regime 1 mean duration & $0.232123$ & $0.241144$ & $0.241582$ \\
\bottomrule
\end{tabular}
\end{center}
both within $\sim\!1$ standard error of target given the sample sizes; the stationary occupation fraction matches almost exactly ($0.2851/0.7149$ empirical vs.\ $0.2857/0.7143$ analytic). For reference, the \emph{old} $\Delta t=0.005$ discretisation's own discrete-to-continuous bias was $0.099999/0.249999$ against the same continuous-time target $0.096458/0.241144$ -- the refined, finer $\Delta t$ is \emph{closer} to the true continuous-time process than the original was, as expected.

\paragraph{Full-pipeline check.} All four policies (oracle, belief-weighted, randomised, naive) run successfully end-to-end on the refined grid; oracle control tables now have shape $(4{,}000, 100)$ per regime (build time $\sim\!13$--$18$s); a single instrumented episode confirms \texttt{episode length == n\_steps == 4{,}000} exactly.

\section{Event-level validation of the adverse-selection markout identity}
\label{sec:eventlevel}

\subsection{Purpose and method}
\texttt{test\_adverse\_selection\_identity.py} (repo root, pytest-compatible) verifies that the simulation's event-level cash, inventory, midprice and reward updates match the adverse-selection markout identity assumed by the theory, $\text{wealth} = C_t + q_tS_t$, and that an adverse fill earns exactly $\delta - \text{jump}$. It builds an isolated single-regime \texttt{TradingEnvironment} (volatility $=0$, \texttt{RunningInventoryPenalty(0,0)} to disable the running/terminal penalties, so the reward reduces exactly to mark-to-market PnL) and forces deterministic arrivals, fills and jump magnitudes by monkey-patching \texttt{get\_arrivals}/\texttt{get\_fills} and replacing the midprice model's \texttt{rng} object with a small mock that returns prescribed values on each call (numpy's \texttt{Generator} is a C-extension type whose bound methods cannot be monkey-patched in place, so the whole object is swapped instead) -- production code is never modified.

\subsection{Results}
Eight cases were checked (ask fill with upward jump; ask arrival/jump without fill; bid fill with downward jump; bid arrival/jump without fill; calm-regime fills without jumps, both sides; existing-inventory revaluation without a fill, both directions; both sides filling in the same step), each asserting the exact expected \texttt{cash}/\texttt{inventory}/\texttt{midprice}/\texttt{wealth\_change} values, plus a reward-reconciliation check (\texttt{reward == wealth\_after - wealth\_before}) in every case. All \textbf{29 checks passed}, with absolute error at floating-point noise level ($\le 1.13\times10^{-14}$).

The most informative case is both sides filling in the same step: with $\delta^a=0.7$, $\delta^b=0.8$, upward jump $0.5$, downward jump $0.4$, the net price jump is $+0.1$, yet
\begin{equation}
\text{wealth\_change} = \delta^a + \delta^b = 1.5 \quad \text{exactly, with no residual jump-related cross term.}
\end{equation}
This holds because inventory returns to exactly zero when both sides fill (an algebraic identity in the fill-multiplier bookkeeping, $\{-1,+1\}$ for $\{$bid, ask$\}$, independent of price), so $\text{wealth}=C_t+q_tS_t$ collapses to $\text{wealth}=C_t$ regardless of the intervening price path.

\subsection{Event ordering}
Confirmed directly from \texttt{TradingEnvironment.\_update\_state}/\texttt{\_update\_agent\_state}/\texttt{\_update\_market\_state} (not merely asserted): (1) arrival generation, (2) fill determination -- both inside \texttt{get\_arrivals\_and\_fills}; (3) cash \& inventory update via \texttt{model\_dynamics.update\_state}, using \texttt{self.midprice} which is still the \emph{pre-jump} value at this point; (4) midprice (diffusion + jump) update, which runs \emph{after} step 3; (5) reward calculation, using \texttt{current\_state} (captured before any of the above) and \texttt{next\_state} (captured after all of the above, so its price already reflects the jump). This ordering is exactly consistent with the HJB's $\delta-\text{jump}$ assumption: a fill locks in execution at the pre-jump price, and only afterward does the price move to reflect the informed order's impact.

\section{Switching-aware oracle: a sensitivity analysis}
\label{sec:switchingaware}

\subsection{Purpose}
\texttt{diagnose\_switching\_oracle.py} is a standalone sensitivity analysis (no production file modified) directly answering the question raised in Section~\ref{sec:discussion}'s ``regime-static, not regime-switching-aware'' discussion: does accounting for future Markov regime switches materially change the optimal quote depths, and could this explain why the oracle does not dominate naive by more?

\subsection{Corrected derivation}
As corrected above (the post-hoc note following \eqref{eq:coupledhjb}), the switching term is additive and linear in $h$, giving the coupled nonlinear system
\begin{equation}
\frac{dh_q^{(i)}}{dt} = \phi q^2 - \frac{r_+^{(i)}}{\kappa}e^{\kappa(h_{q-1}^{(i)}-h_q^{(i)})} - \frac{r_-^{(i)}}{\kappa}e^{\kappa(h_{q+1}^{(i)}-h_q^{(i)})} - \nu_{i,1-i}\big(h_q^{(1-i)}(t) - h_q^{(i)}(t)\big), \quad h_q^{(i)}(T)=-\alpha q^2,
\label{eq:switchinghjb}
\end{equation}
with $\nu_{01},\nu_{10}$ from \eqref{eq:Pnew} and $r_\pm^{(i)}$, the quote formula, and the padded-grid boundary treatment \emph{unchanged} from the static case (the per-instant fill optimisation is regime-local and unaffected by the additive switching term). Solved via backward RK4 time-stepping (not assumed reducible to a linear system in $w$), reusing the existing solver's exact grid/boundary/fill conventions.

\subsection{Validation}
\begin{itemize}
\item \textbf{Zero-switching} ($\nu_{01}=\nu_{10}=0$): reproduces the existing static controls to $3.3\times10^{-4}$ (regime 0) / $2.6\times10^{-4}$ (regime 1) max absolute difference on the interior grid (excluding the outermost $2$ points, which inherit the same padded-boundary conditioning artefact as Error 7 -- the existing solver exponentiates to $w\sim e^{\kappa\phi Q_{\max}^2}$ there, an astronomically large, precision-losing number in $w$-space that this $h$-space solver never produces).
\item \textbf{Identical regimes}: $\max|h_0-h_1| = \max|\delta^a_0-\delta^a_1| = \max|\delta^b_0-\delta^b_1| = 0$ exactly, for any $\nu_{01},\nu_{10}$.
\item \textbf{Terminal condition}: $h_i(T,q)=-\alpha q^2$ exactly, both regimes.
\end{itemize}

\subsection{Main comparison and conclusion}
The switching-aware and static controls differ negligibly near $q=0$ (mean $|{\Delta}|<0.001$) and grow smoothly (not as a boundary artefact) with $|q|$: at $|q|\le 10$, max difference $0.018$ (regime 0) / $0.010$ (regime 1); at $|q|\le 30$, $0.052$/$0.029$. Both thresholds in the pre-registered staged-simulation rule (\texttt{max\_abs\_depth\_difference}$>0.02$, \texttt{max\_abs\_fill\_probability\_diff}$>0.01$) were exceeded (driven by the extreme-inventory region), triggering a paired $200$-episode Monte Carlo comparison (identical exogenous paths, static vs.\ switching-aware oracle):
\begin{equation}
\text{mean } D_{\text{pnl}} = 0.000328\%,\quad \text{mean } D_{\text{objective}} = 0.000327\% \quad \text{(both 95\% CIs include zero).}
\end{equation}
\textbf{Conclusion: switching anticipation is very unlikely to explain the oracle--naive similarity.} The control-table differences are real but confined to inventory levels ($|q|\gtrsim 20$) the simulation essentially never visits (realistic simulated inventory stays within $|q|\approx 2$--$4$), so they never translate into a measurable difference in realised performance.

\section{Oracle-vs-naive paired Monte Carlo comparison}
\label{sec:oraclenaive}

\subsection{Methodology}
\texttt{compare\_oracle\_naive\_paired.py} (results in \texttt{results/oracle\_naive\_paired\_1000.csv}) pairs the oracle and naive policies episode-by-episode so both see \emph{bit-for-bit identical} exogenous randomness -- regime path, arrivals, Brownian shocks, jump sizes, fill uniforms -- differing only in which depths they quote (hence which fills realise). Pairing requires re-seeding \emph{both} the legacy global \texttt{np.random} state (which \texttt{RegimeSwitchingEnv}'s own regime-switching draws use) and the per-seed \texttt{default\_rng} streams (\texttt{make\_regime\_envs(seed=...)}) identically before each policy's fresh environment; this was verified explicitly, not assumed, by instrumenting both runs and diffing the raw regime/arrival/jump/diffusion sequences (0 mismatches across all checked episodes). For $n=1{,}000$ episodes, only the first $20$ carry this full instrumentation (a lighter-weight recorder, wrapping only the jump-magnitude draw needed for \texttt{adverse\_selection\_loss}, is used thereafter for efficiency).

\subsection{Results ($n=1{,}000$)}
\begin{center}
\begin{tabular}{@{}lrrrr@{}}
\toprule
Quantity & mean & std.\ error & 95\% CI \\
\midrule
$D_{\text{pnl}}$ (oracle$-$naive) & $2.629$ & $1.102$ & $[0.469, 4.789]$ \\
$D_{\text{objective}}$ (oracle$-$naive) & $2.812$ & $1.105$ & $[0.647, 4.977]$ \\
\bottomrule
\end{tabular}
\end{center}
Both 95\% CIs exclude zero -- the oracle's advantage on both metrics is statistically significant at $n=1{,}000$ (it was not at $n=500$: $D_{\text{objective}}$ 95\% CI $[-0.31, 6.07]$ included zero there, illustrating how much sampling noise this metric carries per episode).

\subsection{Economic decomposition}
Reconciling exactly (error $0.000\times10^{0}$) against $\text{mean}(D_{\text{objective}})$:
\begin{center}
\begin{tabular}{@{}lr@{}}
\toprule
Component & Value \\
\midrule
Lost spread revenue & $-4.724$ \\
Reduced adverse-selection loss & $+8.001$ \\
Other PnL effects (residual) & $-0.648$ \\
Lower running penalty & $+0.150$ \\
Lower terminal penalty & $+0.033$ \\
\midrule
Net objective advantage & $+2.812$ \\
\bottomrule
\end{tabular}
\end{center}
The oracle's entire net advantage comes from cutting adverse-selection loss almost twice over what it sacrifices in spread revenue; the penalty savings and residual PnL term are minor by comparison. Relative to naive: adverse-selection loss $31.6\%$ lower, total fills $9.4\%$ lower, mean absolute inventory $24.0\%$ lower, terminal absolute inventory $27.0\%$ lower.

\subsection{Units and interpretation notes}
Percentage figures throughout use $\text{NOTIONAL\_CAPITAL} = S_0 \times q_{\max} = \$1{,}000{,}000$ (the value of the \emph{largest possible} position, not capital actually deployed) as a fixed denominator, chosen for comparability across configurations -- not a claim about capital-at-risk. Raw per-episode PnL (oracle $\approx\$42.72$, naive $\approx\$40.09$) is the more directly interpretable quantity. Episodes do \emph{not} compound: each is an independent simulated trading day starting flat (matching the terminal-inventory-aversion $\alpha$ term's ``flatten by close'' convention, Error 6), not a multi-day capital trajectory -- summing/averaging across episodes estimates each policy's \emph{expected per-episode} performance, which is the relevant quantity for policy comparison (a separate compounding analysis would be needed to answer ``how would capital grow over $1{,}000$ consecutive days,'' which additionally requires a position-sizing/leverage assumption not specified here).

\section{Four-policy paired comparison}
\label{sec:fourpolicy}

\subsection{Extension and an RNG-isolation subtlety}
\texttt{compare\_four\_policies\_paired.py} (results in \texttt{results/four\_policy\_paired\_1000.csv}) extends the pairing methodology of Section~\ref{sec:oraclenaive} to all four policies (oracle, belief-weighted, randomised, naive). One subtlety required care: the existing randomised policy samples a hard regime label via \texttt{np.random.choice(2, p=[1-$\pi_t$,$\pi_t$])} -- the \emph{same global} stream \texttt{RegimeSwitchingEnv}'s own regime-switching draws use. Calling it verbatim inside a paired episode would consume extra draws from that shared stream mid-episode, desynchronising the regime path for randomised relative to the other three policies. This script instead draws that sample from a dedicated, independent \texttt{np.random.default\_rng(episode\_seed + 1000)} -- the same pattern already used elsewhere in this codebase (\texttt{plot\_final\_belief.py}) -- which does not change the policy's behaviour (still $\text{Bernoulli}(\pi_t)$ sampling of a hard regime label), only which random-number stream feeds that one decision. Verified: $0$ exogenous-path mismatches across all four policies over the first $20$ (fully-instrumented) episodes.

\subsection{Results ($n=1{,}000$)}
\begin{center}
\begin{tabular}{@{}lrr@{}}
\toprule
Policy & mean raw PnL & mean full objective \\
\midrule
Oracle & $42.722$ & $42.495$ \\
Belief-weighted & $42.437$ & $42.211$ \\
Randomised & $42.333$ & $42.106$ \\
Naive & $40.093$ & $39.683$ \\
\bottomrule
\end{tabular}
\end{center}
Paired differences vs.\ naive, all six 95\% CIs excluding zero:
\begin{center}
\begin{tabular}{@{}lrrr@{}}
\toprule
Policy & $D_{\text{pnl}}$ mean [95\% CI] & $D_{\text{objective}}$ mean [95\% CI] \\
\midrule
Oracle & $2.629\ [0.469, 4.789]$ & $2.812\ [0.647, 4.977]$ \\
Belief-weighted & $2.344\ [0.224, 4.464]$ & $2.528\ [0.403, 4.653]$ \\
Randomised & $2.240\ [0.085, 4.395]$ & $2.423\ [0.264, 4.583]$ \\
\bottomrule
\end{tabular}
\end{center}
The ranking oracle $>$ belief-weighted $>$ randomised $>$ naive holds for both raw PnL and full objective, exactly matching theory (perfect regime knowledge $>$ smooth belief interpolation $>$ hard Bernoulli sampling of the belief $>$ no regime information at all), and matches Section~\ref{sec:discussion}'s ``randomised agent'' discussion (collapsing a continuous belief to a discrete draw discards information relative to smooth interpolation). All three regime-aware policies cut adverse-selection loss and manage inventory almost identically to \emph{each other} (mean $|q|\approx 3.30$ across all three vs.\ naive's $4.34$) -- they are much closer to one another than any is to naive, indicating that \emph{using} the regime/belief signal at all is what matters most; \emph{how} it is used (exact, interpolated, or sampled) is a second-order refinement on top of that.

\section{Epsilon sweep: filter tracking sensitivity to the jump size}
\label{sec:epsilonsweep}

\subsection{Motivation and setup}
A natural question about the corrected filter (Sections~\ref{sec:jumpmix},~\ref{sec:expjump}) is how its tracking quality depends on the size of the adverse-selection jump $\epsilon$: does detection degrade gracefully as the jump signal weakens, and at what scale does it stop mattering? \texttt{diagnostic\_epsilon\_sweep.py} answers this by sweeping $\epsilon$ downward from the calibrated $0.5$ (price units) through six log-spaced values to $0.0005$, building the actual \texttt{mbt\_gym}-driven \texttt{RegimeSwitchingEnv} and the production \texttt{HamiltonFilter} at each value (not a standalone reimplementation), and measuring, over $20$ episodes per $\epsilon$: belief accuracy ($\mathbb{1}[\pi_t\ge0.5]$ vs.\ true regime), belief separation ($\overline{\pi\mid Z{=}1} - \overline{\pi\mid Z{=}0}$), and post-switch detection lag (steps from a $0{\to}1$ regime switch until $\pi_t$ first crosses $0.5$, over $61$ observed switches per $\epsilon$).

This required one small, backward-compatible extension: \texttt{make\_regime\_envs} (\texttt{envs/make\_envs.py}) gained optional \texttt{epsilon}, \texttt{r0\_volatility}, \texttt{r1\_volatility} keyword arguments, each defaulting to the existing module-level constant, so existing call sites are completely unaffected while the sweep script can override them.

\subsection{First pass: calibrated (unequal) volatilities -- flat, uninformative}
\label{sec:epsilonsweep-unequal}
With $\sigma_0^{\text{pct}}=0.01$, $\sigma_1^{\text{pct}}=0.03$ (the calibration in effect at the time of this diagnostic, a $3\times$ ratio -- since superseded, see Section~\ref{sec:equalvolrecal}) held fixed while sweeping $\epsilon$, all three metrics were essentially flat across five decades of $\epsilon$:
\begin{center}
\begin{tabular}{@{}rrrr@{}}
\toprule
$\epsilon$ (price units) & accuracy & separation & mean detection lag (steps) \\
\midrule
$0.500$ & $0.991$ & $0.955$ & $3.1$ \\
$0.126$ & $0.992$ & $0.960$ & $2.0$ \\
$0.032$ & $0.991$ & $0.960$ & $2.4$ \\
$0.008$ & $0.991$ & $0.959$ & $2.9$ \\
$0.002$ & $0.992$ & $0.963$ & $1.7$ \\
$0.0005$ & $0.991$ & $0.957$ & $1.9$ \\
\bottomrule
\end{tabular}
\end{center}
This is not a bug: with $\sigma_1/\sigma_0=3$, the diffusion-only per-step likelihood ratio between regimes is already large enough that the filter locks onto the correct regime within a couple of steps of a switch from the volatility difference alone -- aggregate accuracy is additionally near-saturated by regime persistence (mean dwell time $\sim400$--$1000$ steps; once locked on, hundreds of correlated steps keep the belief there regardless of how weak the per-step signal is). Over this entire range, down to $\epsilon_{\text{pct}}$ on the same order as $\sigma_1^{\text{step}}$ itself, $\epsilon$ contributes no additional identifying information visible in either metric.

\subsection{Second pass: equal volatilities -- isolating epsilon}
\label{sec:epsilonsweep-equal}
Setting both regimes to the same diffusion volatility, $\sigma_0^{\text{pct}}=\sigma_1^{\text{pct}}=0.01$ (so $\sigma_0^{\text{step}}=\sigma_1^{\text{step}}\approx1.58\times10^{-6}$), removes the diffusion channel entirely, leaving $\epsilon$ as the \emph{only} regime-distinguishing signal. Under this ablation, filter performance degrades clearly as $\epsilon$ shrinks:
\begin{center}
\begin{tabular}{@{}rrrr@{}}
\toprule
$\epsilon$ (price units) & accuracy & separation & mean detection lag (steps) \\
\midrule
$0.500$ & $0.946$ & $0.762$ & $9.3$ \\
$0.126$ & $0.949$ & $0.771$ & $10.2$ \\
$0.032$ & $0.947$ & $0.768$ & $11.9$ \\
$0.008$ & $0.941$ & $0.758$ & $14.1$ \\
$0.002$ & $0.937$ & $0.724$ & $12.5$ \\
$0.0005$ & $0.891$ & $0.555$ & $20.0$ \\
\bottomrule
\end{tabular}
\end{center}
Accuracy falls from $0.946$ to $0.891$, separation from $0.762$ to $0.555$, and mean post-switch detection lag roughly doubles ($9.3\to20.0$ steps) as $\epsilon$ drops two orders of magnitude toward $\epsilon_{\text{pct}}\approx\sigma^{\text{step}}$. This confirms the filter's jump-mixture emission (Section~\ref{sec:expjump}) does carry real, epsilon-scaled identifying power on its own -- Section~\ref{sec:epsilonsweep-unequal}'s flat result was a consequence of the diffusion channel dominating at that diagnostic's $3\times$ volatility ratio, not of the jump signal being uninformative in general. This result directly motivated the permanent recalibration in Section~\ref{sec:equalvolrecal}.

\section{Epsilon sweep: policy performance sensitivity to the jump size}
\label{sec:policyepsilonsweep}

\subsection{Motivation and setup}
\texttt{simulate\_policies\_epsilon\_sweep.py} extends the epsilon sweep of Section~\ref{sec:epsilonsweep} from filter tracking quality to realised policy performance: does each policy's edge over naive (Section~\ref{sec:fourpolicy}) shrink as the adverse-selection signal weakens? Like Section~\ref{sec:epsilonsweep-equal}, both regimes share the same diffusion volatility, so $\epsilon$ is the only channel through which regime-aware quoting (oracle, belief-weighted, randomised) can outperform naive.

The script is standalone and does not modify \texttt{simulate\_belief\_weighted.py} or \texttt{compare\_four\_policies\_paired.py}: it imports their pure functions (\texttt{build\_optimal\_control}, \texttt{get\_control}, \texttt{normalise\_depth}) unchanged, and rebuilds only the $\epsilon$-dependent objects locally per sweep value -- the regime-$1$ control table (regime-$0$'s is $\epsilon$-independent and built once, reused across the sweep), the naive depth $1/\kappa+\epsilon/2$, and the Hamilton filter's jump size -- since \texttt{simulate\_belief\_weighted.py} hardcodes \texttt{EPSILON} from \texttt{envs/make\_envs.py} at import time. It reuses \texttt{compare\_four\_policies\_paired.py}'s paired-episode methodology (Section~\ref{sec:fourpolicy}): for each episode index, all four policies see bit-for-bit identical exogenous draws, differing only in quoted depths. Six log-spaced $\epsilon$ values ($0.5\to0.0005$, matching Section~\ref{sec:epsilonsweep}) $\times$ $20$ paired episodes $\times$ $4$ policies were run.

\subsection{Results}
Paired mean objective advantage over naive (\% of notional capital), $20$ episodes per $\epsilon$:
\begin{center}
\begin{tabular}{@{}rrrr@{}}
\toprule
$\epsilon$ (price units) & Oracle & Belief-weighted & Randomised \\
\midrule
$0.5000$ & $0.001053\ [0.000010,\,0.002095]$ & $0.000599\ [-0.000509,\,0.001706]$ & $0.001008\ [-0.000173,\,0.002188]$ \\
$0.1256$ & $0.000094\ [-0.000100,\,0.000288]$ & $0.000082\ [-0.000071,\,0.000236]$ & $0.000101\ [-0.000086,\,0.000287]$ \\
$0.0315$ & $0.000021\ [-0.000079,\,0.000122]$ & $0.000029\ [-0.000068,\,0.000125]$ & $0.000031\ [-0.000069,\,0.000131]$ \\
$0.0079$ & $-0.000021\ [-0.000110,\,0.000067]$ & $-0.000019\ [-0.000107,\,0.000070]$ & $-0.000018\ [-0.000105,\,0.000069]$ \\
$0.0020$ & $-0.000055\ [-0.000150,\,0.000041]$ & $-0.000055\ [-0.000150,\,0.000041]$ & $-0.000055\ [-0.000150,\,0.000041]$ \\
$0.0005$ & $-0.000025\ [-0.000118,\,0.000068]$ & $-0.000025\ [-0.000118,\,0.000068]$ & $-0.000025\ [-0.000118,\,0.000068]$ \\
\bottomrule
\end{tabular}
\end{center}
At the calibrated jump size ($\epsilon=0.5$), only the oracle's $95\%$ CI excludes zero at this sample size ($n=20$; cf.\ Section~\ref{sec:oraclenaive}'s observation that the oracle-naive gap was not significant at $n=500$ either, illustrating how much per-episode noise this metric carries), but all three regime-aware policies' point estimates sit clearly above naive. By $\epsilon\approx0.03$ the advantage has collapsed to within noise of zero, and it stays there (even turning marginally negative, well within the CI) for every smaller $\epsilon$ tested.

\subsection{Interpretation}
This is the expected signature under the equal-volatility ablation: with $\sigma_0=\sigma_1$, a regime-aware policy's only channel for beating naive is trading off spread revenue against avoided adverse-selection losses via the $\epsilon$ term in the optimal depth, $\delta^{a,*}=1/\kappa+\epsilon+\ldots$ (Section~\ref{sec:control}), and via preferentially avoiding fills immediately before a jump. As $\epsilon\to0$, both the optimal control's own $\epsilon$-adjustment and the realised adverse-selection loss shrink to zero for \emph{every} policy alike, so there is decreasingly little for regime-awareness to protect against -- naive's fixed spread $1/\kappa+\epsilon/2$ converges to the same quote every other policy converges to, and the four policies become interchangeable. Combined with Section~\ref{sec:epsilonsweep-equal}'s finding that the filter itself degrades (but does not collapse) over the same range, this indicates the vanishing edge here is driven primarily by the shrinking \emph{economic} value of regime information (there is less adverse selection to avoid), not solely by the filter's detection quality falling off.

See \texttt{results/policies\_epsilon\_sweep.csv} for the full episode-level data and \texttt{policies\_epsilon\_sweep\_*.png} for the accompanying figure (mean risk-adjusted objective vs.\ $\epsilon$ per policy; paired advantage over naive vs.\ $\epsilon$ with $95\%$ CI bands).

\section{Recalibration: equal regime volatilities}
\label{sec:equalvolrecal}

\subsection{Decision}
Following the diagnostics of Sections~\ref{sec:epsilonsweep}--\ref{sec:policyepsilonsweep}, \texttt{envs/make\_envs.py}'s \texttt{R1\_VOLATILITY} was changed from $0.03$ (a $3\times$ ratio against $\sigma_0=0.01$) to $0.01$, equal to \texttt{R0\_VOLATILITY}. \texttt{EPSILON} remains at its existing calibrated value, $0.5$. This is a \emph{permanent} change to the live default, not a one-off ablation script override (unlike Sections~\ref{sec:epsilonsweep-equal},~\ref{sec:policyepsilonsweep}, which forced equal volatilities locally via \texttt{make\_regime\_envs}'s \texttt{r0\_volatility}/\texttt{r1\_volatility} keyword arguments without touching the module default) -- every script that imports \texttt{R1\_VOLATILITY} or \texttt{R1\_SIGMA\_STEP} from \texttt{envs/make\_envs.py} (i.e.\ essentially every script in the project) now sees the equal-volatility environment.

\subsection{Rationale}
Section~\ref{sec:epsilonsweep-equal}'s ablation established that the Hamilton filter can learn the regime from the jump signal alone -- accuracy $0.89$--$0.95$ across a wide range of $\epsilon$ with no diffusion-volatility difference to lean on at all. Given that, retaining an independent $3\times$ volatility difference between regimes confounds two distinct adverse-selection signals (a jump process and an unrelated diffusion-intensity change) into a single ``regime,'' when the substantive object of interest in this dissertation is specifically adverse selection via price impact on informed order flow (\eqref{eq:midprice}'s jump term), not a generic volatility-regime-switching model. Making $\epsilon$ the \emph{only} parameter that varies between regimes is a cleaner and more defensible identification strategy: any measured effect of ``regime awareness'' (filter tracking quality, Section~\ref{sec:epsilonsweep}; policy performance, Section~\ref{sec:policyepsilonsweep}; the four-policy and oracle-naive comparisons of Sections~\ref{sec:oraclenaive}--\ref{sec:fourpolicy}) can now be attributed unambiguously to adverse-selection detection and avoidance, rather than partly to an unmodelled volatility-forecasting effect riding along with it.

\subsection{Consequence: earlier numerical results are now superseded}
\label{sec:equalvolrecal-stale}
All quantitative results reported in Sections~\ref{sec:norm}--\ref{sec:discretisation},~\ref{sec:eventlevel}--\ref{sec:fourpolicy}, and~\ref{sec:epsilonsweep-unequal} were computed under the \emph{previous} calibration ($\sigma_1^{\text{pct}}=0.03$) and are retained here as a historical record of the bug-fix and validation process, not as a description of the current environment's behaviour. In particular, the four-policy paired comparison (Section~\ref{sec:fourpolicy}), the oracle-naive decomposition (Section~\ref{sec:oraclenaive}), and the switching-aware sensitivity analysis (Section~\ref{sec:switchingaware}) have not been re-run under the new equal-volatility calibration; their qualitative conclusions (ranking oracle $>$ belief-weighted $>$ randomised $>$ naive; the switching-aware/static control difference being negligible at realistic inventory levels) are plausible but unverified under the current default, since -- per Section~\ref{sec:policyepsilonsweep} -- the size of each policy's edge over naive is itself sensitive to how much of a diffusion-volatility difference remains between regimes. Sections~\ref{sec:epsilonsweep-equal} and~\ref{sec:policyepsilonsweep} are unaffected by this caveat, since both already used the equal-volatility configuration that is now the live default.

\section{Diagnostic scripts (this session)}
\label{sec:diagnosticscripts}
None of the following modify any production file; each is a standalone, read-only-import diagnostic or sensitivity analysis:
\begin{itemize}
\item \texttt{test\_adverse\_selection\_identity.py} -- Section~\ref{sec:eventlevel}.
\item \texttt{diagnose\_switching\_oracle.py} -- Section~\ref{sec:switchingaware}.
\item \texttt{compare\_oracle\_naive\_paired.py} -- Section~\ref{sec:oraclenaive} (CLI: \texttt{--episodes}, \texttt{--verify-episodes}, \texttt{--base-seed}, \texttt{--output}).
\item \texttt{compare\_four\_policies\_paired.py} -- Section~\ref{sec:fourpolicy}.
\item \texttt{plot\_oracle\_naive\_results.py} -- post-processes \texttt{results/oracle\_naive\_paired\_1000.csv} only (no environment/policy/filter/control imports) into the CI, waterfall-decomposition and relative-activity figures cited in Section~\ref{sec:oraclenaive}.
\item \texttt{diagnostic\_epsilon\_sweep.py} -- Section~\ref{sec:epsilonsweep} (CLI: \texttt{--episodes}).
\item \texttt{simulate\_policies\_epsilon\_sweep.py} -- Section~\ref{sec:policyepsilonsweep} (CLI: \texttt{--episodes}, \texttt{--base-seed}, \texttt{--output}).
\end{itemize}

\end{document}
